\chapter*{Résumé}
\addcontentsline{toc}{chapter}{Résumé}

La gestion quotidienne d'une agence bancaire implique le traitement de données clients sensibles, la tenue rigoureuse des comptes et l'exécution d'opérations financières soumises à des règles métier strictes. Les méthodes traditionnelles, souvent manuelles ou dispersées, augmentent les risques d'erreurs de saisie, d'incohérences de solde et de faible traçabilité.

Face à ces limites, ce projet propose \textbf{\nomapplication}, une application web centralisée développée avec \textbf{Java} et \textbf{Spring Boot} par \etudianteun\ et \etudiantedeux\ dans le cadre d'un projet de fin d'année à l'\ecole.

La solution couvre la gestion interne de l'agence (clients, comptes, transactions, utilisateurs, audit, reporting) ainsi que des extensions : paiement de factures, commande de chéquier, centre de notifications et portail client en ligne. Chaque acteur dispose d'un espace adapté à son rôle.

Le présent rapport présente le contexte du projet, l'analyse des besoins, la conception UML, l'environnement technique, l'architecture logicielle, la réalisation des interfaces et les perspectives d'évolution.

\vspace{0.5cm}
\textbf{Mots-clés :} Agence bancaire — Spring Boot — Gestion de comptes — Transactions — Portail client — UML
